% ============================================================
%  Code Appendix — FR Power Thesis
%  The Role of Weather in French Day-Ahead Electricity Price Forecasting
%  Leo Camberleng & Lyam Oumedjeber — EDHEC MSc Data Analysis & AI
% ============================================================

\documentclass[11pt, a4paper, oneside]{report}

% ── Encoding & language ──────────────────────────────────────
\usepackage[T1]{fontenc}
\usepackage[utf8]{inputenc}
\usepackage[english]{babel}

% ── Typography ───────────────────────────────────────────────
\usepackage{lmodern}
\usepackage{setspace}
\onehalfspacing

% ── Page layout ──────────────────────────────────────────────
\usepackage[
  top=2.5cm, bottom=2.5cm,
  left=3.0cm, right=2.5cm
]{geometry}

% ── Headers & footers ────────────────────────────────────────
\usepackage{fancyhdr}
\pagestyle{fancy}
\fancyhf{}
\fancyhead[L]{\small\textit{Code Appendix — FR Power Thesis}}
\fancyhead[R]{\small\textit{Camberleng \& Oumedjeber}}
\fancyfoot[C]{\thepage}
\renewcommand{\headrulewidth}{0.4pt}

% ── Code listings ────────────────────────────────────────────
\usepackage{listings}
\usepackage[dvipsnames]{xcolor}

\definecolor{codebg}{RGB}{248, 248, 248}
\definecolor{codecomment}{RGB}{100, 100, 100}
\definecolor{codestring}{RGB}{163, 21, 21}
\definecolor{codekw}{RGB}{0, 0, 200}

\lstset{
  language=Python,
  basicstyle=\ttfamily\small,
  keywordstyle=\color{codekw}\bfseries,
  stringstyle=\color{codestring},
  commentstyle=\color{codecomment}\itshape,
  numberstyle=\tiny\color{gray},
  numbers=left,
  stepnumber=5,
  numbersep=8pt,
  backgroundcolor=\color{codebg},
  frame=single,
  framesep=4pt,
  rulecolor=\color{gray!40},
  breaklines=true,
  breakatwhitespace=false,
  showstringspaces=false,
  tabsize=4,
  captionpos=b,
  belowcaptionskip=8pt,
  abovecaptionskip=4pt,
  xleftmargin=12pt,
  xrightmargin=4pt,
}

% ── Hyperlinks ───────────────────────────────────────────────
\usepackage[
  colorlinks=true,
  linkcolor=black,
  urlcolor=black,
  bookmarks=true
]{hyperref}

% ── Misc ─────────────────────────────────────────────────────
\usepackage{enumitem}
\usepackage{booktabs}
\usepackage{float}

% ============================================================
\begin{document}

% ── Title page ───────────────────────────────────────────────
\begin{titlepage}
  \centering
  \vspace*{2cm}

  {\rule{\linewidth}{2pt}}
  \vspace{0.5cm}

  {\LARGE\bfseries
    The Role of Weather in French Day-Ahead\\[6pt]
    Electricity Price Forecasting:\\[6pt]
    A Random Forest Approach\\[12pt]
    \Large --- Code Appendix ---
  \par}

  \vspace{0.5cm}
  {\rule{\linewidth}{1pt}}
  \vspace{2cm}

  {\large Master's Thesis --- MSc Data Analysis \& AI\\[4pt]
   EDHEC Business School}

  \vspace{2cm}

  {\large\bfseries Leo Camberleng}\\[4pt]
  {\large\bfseries Lyam Oumedjeber}

  \vspace{2cm}

  {\large May 2026}

  \vfill
  {\rule{\linewidth}{2pt}}
\end{titlepage}

% ── TOC ──────────────────────────────────────────────────────
\tableofcontents
\clearpage

% ── Preamble ─────────────────────────────────────────────────
\chapter*{Overview}
\addcontentsline{toc}{chapter}{Overview}

This document contains the complete Python codebase used to produce the results
of the thesis \textit{``The Role of Weather in French Day-Ahead Electricity Price
Forecasting: A Random Forest Approach''}.

\medskip
\noindent The codebase is organised as follows:

\begin{itemize}
  \item \textbf{Configuration} (\texttt{src/config.py}): global paths, API keys, and model hyperparameters.
  \item \textbf{Data Collection} (\texttt{scripts/fetch\_*.py}): retrieval from ENTSO-E API, ERA5 CDS API, and fuel price sources.
  \item \textbf{Feature Engineering} (\texttt{src/features.py}, \texttt{src/preprocessing.py}): construction of the 35-feature matrix.
  \item \textbf{Modelling \& Evaluation} (\texttt{src/models.py}, \texttt{src/evaluate.py}): Random Forest and XGBoost training, metrics, Diebold-Mariano tests.
  \item \textbf{Trading Backtest} (\texttt{scripts/backtest.py}): day-ahead directional strategy, P\&L, Sharpe, drawdown, transaction cost scenarios.
  \item \textbf{2022 Crisis Validation} (\texttt{scripts/test\_2022.py}): regime-separated evaluation on the 2022 energy crisis.
  \item \textbf{Pipeline} (\texttt{scripts/run\_full\_pipeline.py}): end-to-end orchestration script.
\end{itemize}

\medskip
\noindent All data sources (ENTSO-E Transparency Platform, ERA5 reanalysis,
fuel prices) are publicly available. The feature matrix is reproducible by
running \texttt{scripts/run\_full\_pipeline.py} after configuring the API keys.

\clearpage

% ============================================================
\chapter{Configuration}
% ============================================================

\section{\texttt{src/config.py}}

Global paths, API credentials (loaded from environment variables), market
parameters, and model hyperparameters. Random Forest uses 500 trees with
\texttt{max\_depth=10} and \texttt{min\_samples\_leaf=5}. XGBoost uses 500
estimators with learning rate $\eta=0.05$ and \texttt{max\_depth=6}.

\lstinputlisting[
  caption={\texttt{src/config.py} — Project configuration and hyperparameters},
  label={lst:config}
]{../src/config.py}

\clearpage

% ============================================================
\chapter{Data Collection}
% ============================================================

\section{\texttt{scripts/fetch\_entsoe.py}}

Retrieves hourly day-ahead prices, load forecasts, actual generation by source,
and cross-border physical flows for France (bidding zone
\texttt{10YFR-RTE------C}) from the ENTSO-E Transparency Platform using the
\texttt{entsoe-py} Python library.

\lstinputlisting[
  caption={\texttt{scripts/fetch\_entsoe.py} — ENTSO-E data retrieval},
  label={lst:fetch_entsoe}
]{../scripts/fetch_entsoe.py}

\clearpage

\section{\texttt{scripts/fetch\_era5.py}}

Downloads ERA5 reanalysis weather data (2m temperature, 10m wind speed,
solar irradiance, precipitation) for metropolitan France from the Copernicus
Climate Data Store using the \texttt{cdsapi} library.

\lstinputlisting[
  caption={\texttt{scripts/fetch\_era5.py} — ERA5 weather data download},
  label={lst:fetch_era5}
]{../scripts/fetch_era5.py}

\clearpage

\section{\texttt{scripts/fetch\_fuels.py}}

Retrieves TTF natural gas and ARA coal price series from public market data
sources and aligns them to the hourly electricity price index.

\lstinputlisting[
  caption={\texttt{scripts/fetch\_fuels.py} — Fuel price data retrieval},
  label={lst:fetch_fuels}
]{../scripts/fetch_fuels.py}

\clearpage

% ============================================================
\chapter{Feature Engineering}
% ============================================================

\section{\texttt{src/preprocessing.py}}

Data cleaning, timestamp alignment, and interpolation of missing values.
All series are aligned to UTC+1 (Europe/Paris) and reindexed to a complete
hourly grid.

\lstinputlisting[
  caption={\texttt{src/preprocessing.py} — Data cleaning and alignment},
  label={lst:preprocessing}
]{../src/preprocessing.py}

\clearpage

\section{\texttt{src/features.py}}

Constructs the full 35-feature matrix from the cleaned data sources.
Key features include: calendar encodings (sine/cosine), price lags
(24h, 48h, 168h), rolling means, nuclear availability ratio,
Heating Degree Days (HDD, threshold 17\textdegree C), Cooling Degree Days
(CDD, threshold 22\textdegree C), wind power proxy ($v^3$), and the
Weather Stress Index (WSI = HDD / wind speed).

\lstinputlisting[
  caption={\texttt{src/features.py} — Feature matrix construction (35 features)},
  label={lst:features}
]{../src/features.py}

\clearpage

\section{\texttt{scripts/build\_features.py}}

Orchestration script that calls the preprocessing and feature engineering
modules to produce the final \texttt{features.parquet} file.

\lstinputlisting[
  caption={\texttt{scripts/build\_features.py} — Feature pipeline orchestration},
  label={lst:build_features}
]{../scripts/build_features.py}

\clearpage

% ============================================================
\chapter{Modelling and Evaluation}
% ============================================================

\section{\texttt{src/models.py}}

Wraps scikit-learn's \texttt{RandomForestRegressor} and XGBoost's
\texttt{XGBRegressor} with a common interface. Handles the strict
temporal train/test split (no shuffling).

\lstinputlisting[
  caption={\texttt{src/models.py} — Model training wrappers},
  label={lst:models}
]{../src/models.py}

\clearpage

\section{\texttt{src/evaluate.py}}

Computes MAE, RMSE, sMAPE, R$^2$, and hit ratio on the test set.
Implements the Diebold-Mariano test with Harvey-Leybourne-Newbold
correction for pairwise model comparison.

\lstinputlisting[
  caption={\texttt{src/evaluate.py} — Forecast evaluation metrics and DM test},
  label={lst:evaluate}
]{../src/evaluate.py}

\clearpage

\section{\texttt{src/plots.py}}

Generates all thesis figures: forecast vs.\ actual, error distributions,
monthly MAE, feature importance, and weather stress index.

\lstinputlisting[
  caption={\texttt{src/plots.py} — Figure generation},
  label={lst:plots}
]{../src/plots.py}

\clearpage

% ============================================================
\chapter{Trading Backtest}
% ============================================================

\section{\texttt{scripts/backtest.py}}

Implements the day-ahead directional trading strategy described in
Chapter~5 of the thesis. Key design choices:
\begin{itemize}
  \item Signal: $\text{sign}(\hat{p}(h,D+1) - p(h,D))$ — executable before
        12:00 CET gate closure
  \item Dead-band filter: $|\hat{p}(h,D+1) - p(h,D)| < 2$ EUR/MWh $\Rightarrow$ flat
  \item P\&L: $\text{signal} \times (p(h,D+1) - p(h,D))$, units EUR/MW
  \item Transaction costs: 0.10 / 0.30 / 0.60 EUR/MWh per active hour
  \item Sharpe: daily aggregation, annualised $\times\sqrt{365}$
\end{itemize}

\lstinputlisting[
  caption={\texttt{scripts/backtest.py} — Day-ahead directional trading backtest},
  label={lst:backtest}
]{../scripts/backtest.py}

\clearpage

% ============================================================
\chapter{2022 Energy Crisis Validation}
% ============================================================

\section{\texttt{scripts/test\_2022.py}}

Retrains Models A--D on 2018--2021 data and evaluates them on the full
year 2022 (8,760 hours). Produces the performance table and DM tests
for the 2022 crisis regime reported in Section~4.7 of the thesis.
The key result: weather features are highly significant in 2022
(DM $= -13.27$, $p < 0.001$) versus non-significant in the stable
2024--25 period ($p = 0.572$).

\lstinputlisting[
  caption={\texttt{scripts/test\_2022.py} — 2022 energy crisis out-of-sample evaluation},
  label={lst:test2022}
]{../scripts/test_2022.py}

\clearpage

% ============================================================
\chapter{End-to-End Pipeline}
% ============================================================

\section{\texttt{scripts/run\_full\_pipeline.py}}

Master orchestration script that runs the complete pipeline in sequence:
data loading, feature engineering, model training (Models A--D),
evaluation, figure generation, and backtest. Running this script from
the project root reproduces all results in the thesis.

\lstinputlisting[
  caption={\texttt{scripts/run\_full\_pipeline.py} — End-to-end pipeline},
  label={lst:pipeline}
]{../scripts/run_full_pipeline.py}

\end{document}
